\subsection{Design principles}
\label{sec:principles}

Well-formedness enforces several constraints, summarised here with the programs each accepts and rejects;
the rules follow in \secref{wf-patterns} to \secref{wf-modules}. Two auxiliary functions appear
throughout: $\assignsS(s)$, the set of variables assigned anywhere in $s$, not descending into nested
function definitions, which have their own scope, so that $\assignsS(s) = \dom(\Delta)$ when
$s : \Assigns{\Delta}$; and $\binds(p)$, the set of variables a pattern $p$ introduces.

\subsubsection{No reassignment of captured variables}
\label{sec:captured-vars}

A lambda or nested function definition may refer to variables of the enclosing scope, which it is said to
\emph{capture}. The set of variables of the enclosing scope captured by the closures within a statement
$s$ is $\capturesS(s)$, with $\capturesE(e)$ the same for an expression $e$: the free variables of each
closure's body, less its parameters and the variables the body assigns, which are local to it. For a
mutual region of function definitions the names of the functions themselves are also excluded, since the
\ruleName{def} rule binds them for every body in the region (\secref{mutual-recursion}).

In Python a closure captures a mutable binding, so reassigning a captured variable after the closure is
created changes what the closure sees, as in the example of \secref{no-mutable-state}, where \pyinline{x}
is captured by the definition of \pyinline{g} and then reassigned by \pyinline{x = 6}. \PurePy rejects
that program by a side condition on the \ruleName{cons} rule: in a block $s; b$, no variable captured by
$s$ may be assigned in $b$, that is $\capturesS(s) \cap \assignsS(b) = \varnothing$. The
\ruleName{statement} rule imposes the same condition on the statements of a module. A dataclass
declaration counts as an assignment of its class name here, though re-declaring a class name is excluded
outright (\secref{no-first-class}).

The same issue arises within a single statement when the expression assigned to a variable captures that
variable:

\begin{python}
def f():
    x = 5
    x = lambda: x  # Python: the lambda returns itself
    return x()
\end{python}

\noindent The \ruleName{assign} rule requires $x \notin \capturesE(e)$ for an assignment $x = e$, so the
program is rejected. A function definition is not subject to this condition, because the \ruleName{def}
rule binds the names of its mutual region for every body: a body's reference to its own name or a
sibling's denotes the definition, which is also what Python's late binding finds
(\secref{mutual-recursion}).
